% 作品介绍：面向线性PCurve的一类通用样条面到二次曲面智能转化方案
% 第一章：2026-10-16 用户确认稿。此文件为章节片段，待整稿组装。
\section{项目背景}

在 CAD 模型的数据交换与导入过程中，部分具有解析几何结构的曲面以非均匀有理 B 样条（NURBS）形式存储。这些曲面可能来源于解析曲面的样条化转换、特定构型下的圆角过渡，或采用样条支撑曲面的裁剪面。尽管其几何形状仍与平面、圆柱面、圆锥面或球面等解析曲面一致，或在给定容差内接近，原有的曲面类型与几何参数却未必得到显式保留。

恢复此类面片的解析表示具有直接的工程价值。一方面，解析曲面通常可由少量几何参数描述，有助于减少控制顶点、节点和权重等数据的存储开销；另一方面，明确的曲面类型能够为求交、投影及距离计算提供专用求解路径，减少部分通用数值计算的开销，为后续几何处理提供更有效的基础表示。

然而，支撑曲面的近似吻合并不足以保证有限面片转换有效。转换过程还需处理边界对应、参数化差异及数值退化，并在规定容差内完成相应的几何验收。为此，本项目面向具有线性 PCurve 结构的一类样条面片，构建集解析曲面识别、几何参数恢复、必要优化与独立验收于一体的自动转化方案，为导入模型的几何表示简化及后续算法加速提供技术支撑。
